\documentclass[11pt,a4paper]{article}
\usepackage[margin=1.1in]{geometry}
\usepackage{amsmath,amssymb,amsthm,mathtools}
\usepackage[T1]{fontenc}
\usepackage{microtype}
\usepackage{enumitem}
\usepackage{booktabs}
\usepackage[hidelinks]{hyperref}

\newtheorem{theorem}{Theorem}[section]
\newtheorem{proposition}[theorem]{Proposition}
\newtheorem{lemma}[theorem]{Lemma}
\newtheorem{corollary}[theorem]{Corollary}
\theoremstyle{definition}
\newtheorem{definition}[theorem]{Definition}
\newtheorem{remark}[theorem]{Remark}
\theoremstyle{remark}
\newtheorem{scope}[theorem]{Scope}
\newtheorem{residue}{Residue}

\newcommand{\Z}{\mathbb{Z}}
\newcommand{\Q}{\mathbb{Q}}
\newcommand{\C}{\mathbb{C}}
\newcommand{\R}{\mathbb{R}}
\newcommand{\HH}{\mathbb{H}}
\newcommand{\vol}{\operatorname{vol}}
\newcommand{\Sym}{\operatorname{Sym}}
\newcommand{\Rz}{\mathcal{R}}

\title{The boundary-graviton one-loop determinant\\
as an infinite product of twisted Ruelle zetas,\\
and the three residues that remain}
\author{}
\date{26 August 2026}

\begin{document}
\maketitle

\begin{abstract}
For a hyperbolic $3$-manifold $X=\Gamma\backslash\HH^3$, the one-loop determinant of boundary
gravitons in $\mathrm{AdS}_3$ is conventionally written as a product over primitive conjugacy
classes in terms of the nome $q_\gamma=e^{-\ell_\gamma+i\theta_\gamma}$. We observe that the
symmetric-power representation $\sigma_k$ entering the twisted Ruelle zeta at integer argument
$k$ is \emph{one-dimensional}, so that
\[
   \Rz(k,\sigma_k)=\prod_{[\gamma]\ \mathrm{prime}}\bigl(1-q_\gamma^{\,k}\bigr)
\]
with $q_\gamma$ \emph{exactly} that nome. The boundary-graviton one-loop determinant is
therefore
\[
   Z_{\mathrm{1-loop}}=\prod_{n\ge2}\bigl|\Rz(n,\sigma_n)\bigr|^{-2},
\]
an infinite product of twisted Ruelle zeta values. In particular it is \emph{not} any single
Ray--Singer analytic torsion: to the question ``which of Pfaff's torsions is it?'' the answer is
none of them. Pfaff's ratio formula supplies exactly the tail $k\ge3$; the one factor it cannot
reach is $n=2$, which sits \emph{at} the abscissa of absolute convergence $\mathrm{Re}(s)>2$ ---
which is also why that formula begins at $m\ge3$ and normalises by $\rho(2)$.

We then state precisely what is not thereby established. Three residues remain, and we
separate them rather than merging them into a single caveat: (1) a gauge-fixed spin-$2$ test
function for the finite-volume cusped case, which we believe is a genuine gap in the
literature; (2) an evaluation-point gap --- Fried's theorem relates the twisted Ruelle zeta at
$s=0$ to analytic torsion, Pfaff's ratios live at $k\ge3$, and the graviton product at
$s=n\ge2$, so these are three evaluations of one family and Fried's is the one point the
graviton never visits; (3) the cusp's continuous spectrum. For $m004$ we also settle a related
question negatively: its Ruelle zeta does not factor through $L(s,\chi_{-3})$ in its Euler
product, which is purely geodesic, although $L(s,\chi_{-3})$ does enter the functional equation
through the cusp. No measured quantity is claimed anywhere.
\end{abstract}

\section{Introduction}

\subsection{Setting}

Let $X=\Gamma\backslash\HH^3$ be a hyperbolic $3$-manifold. For a primitive closed geodesic
$[\gamma]$ of complex length $\ell_\gamma+i\theta_\gamma$, write
$q_\gamma=e^{-\ell_\gamma+i\theta_\gamma}$. In the standard treatment of one-loop quantum
gravity in $\mathrm{AdS}_3$, the contribution of boundary gravitons is assembled from products
over $[\gamma]$ in powers of $q_\gamma$ and $\bar q_\gamma$.

Separately, for a finite-dimensional representation $\sigma$ of $\Gamma$ one has the twisted
Ruelle zeta function, defined for $\mathrm{Re}(s)$ large by
\[
   \Rz(s,\sigma)=\prod_{[\gamma]\ \mathrm{prime}}
      \det\bigl(I-\sigma(\gamma)\,e^{-s\,\ell_\gamma}\bigr),
\]
and continued meromorphically. Ruelle zeta values are related to Ray--Singer analytic torsion
by theorems of Fried and, in the setting relevant here, by Pfaff's ratio formula.

These two objects are usually discussed in different languages. The point of \S\ref{sec:identity}
is that in this case they are not merely related --- they coincide, term by term, for an
elementary reason.

\subsection{The observation}

The representations $\sigma_k$ that appear are one-dimensional. When $\sigma$ is
one-dimensional the determinant in the definition of $\Rz$ is just a scalar, and the
$e^{-s\ell_\gamma}$ combines with the character value $e^{ik\theta_\gamma}$ into precisely
$q_\gamma^{\,k}$. There is no room for the two conventions to differ, and the identification is
an identity rather than a correspondence. Everything in \S\ref{sec:identity} follows from that
observation; we claim no depth for it, only that stating it clearly changes what one can then
prove and, more usefully, changes what one can then \emph{locate} as missing.

\subsection{What this paper claims, and what it does not}

\begin{itemize}[leftmargin=1.4em]
\item \textbf{Claimed:} the identity (Theorem~\ref{thm:identity}); the consequence that the
      one-loop determinant is an infinite product and not a single torsion
      (Corollary~\ref{cor:notasingle}); the exact match between Pfaff's ratio formula and the
      $k\ge3$ tail (Proposition~\ref{prop:tail}); numerical evidence, on an instrument with a
      control proved able to detect failure, that the $n=2$ factor does not break down at the
      abscissa (\S\ref{sec:n2}); and the negative result on $L(s,\chi_{-3})$
      (Theorem~\ref{thm:Lchi}).
\item \textbf{Not claimed:} the assembly of the one-loop partition function. Three specific
      obstructions remain and are stated as Residues~\ref{res:one}--\ref{res:three} in
      \S\ref{sec:residues}, each with what would close it. Convergence of the $n=2$ factor is
      \emph{not} proved (Scope~\ref{sc:n2}). No numerical value is offered as a measurement of
      anything.
\end{itemize}

\section{The identity}\label{sec:identity}

\begin{theorem}[the dictionary entry is an identity]\label{thm:identity}
Let $\sigma_k$ denote the one-dimensional representation sending $\gamma\mapsto
e^{ik\theta_\gamma}$. Then for $\mathrm{Re}(s)>2$,
\[
   \Rz(k,\sigma_k)=\prod_{[\gamma]\ \mathrm{prime}}\bigl(1-q_\gamma^{\,k}\bigr),
   \qquad q_\gamma=e^{-\ell_\gamma+i\theta_\gamma},
\]
with $q_\gamma$ exactly the nome appearing in the boundary-graviton product. Consequently
\[
   Z_{\mathrm{1-loop}}=\prod_{n\ge2}\bigl|\Rz(n,\sigma_n)\bigr|^{-2}.
\]
\end{theorem}

\begin{proof}
Since $\sigma_k$ is one-dimensional, $\det(I-\sigma_k(\gamma)e^{-k\ell_\gamma})
=1-e^{ik\theta_\gamma}e^{-k\ell_\gamma}=1-\bigl(e^{-\ell_\gamma+i\theta_\gamma}\bigr)^k
=1-q_\gamma^{\,k}$. The boundary-graviton determinant is the product over $n\ge2$ of
$\prod_\gamma|1-q_\gamma^{\,n}|^{-2}$, which is the displayed expression.
\end{proof}

\begin{corollary}[it is not a single analytic torsion]\label{cor:notasingle}
$Z_{\mathrm{1-loop}}$ is an infinite product of Ruelle zeta values. It is therefore not equal to
any one Ray--Singer analytic torsion $T_X(\rho)$, nor to any finite ratio of such; asking which
torsion it ``is'' has the answer: none of them.
\end{corollary}

This is a negative answer to a natural question, and it is worth stating because the question is
asked. The positive content is that the individual factors are torsion-theoretic, which is what
Proposition~\ref{prop:tail} makes precise.

\begin{proposition}[Pfaff's formula supplies exactly the tail]\label{prop:tail}
With $\rho(m)=\Sym^{2m}\C^2$ and $\kappa(X)$ the number of cusps, Pfaff's ratio formula reads
\[
   \frac{T_X(\rho(m))}{T_X(\rho(2))}
   =\Bigl(\frac{c(m)}{c(2)}\Bigr)^{\kappa(X)}
     \exp\!\Bigl(-\tfrac1\pi\vol(X)\bigl(m(m+1)-6\bigr)\Bigr)
     \prod_{k=3}^{m}\bigl|\Rz(k,\sigma_k)\bigr| .
\]
The product on the right is exactly the $k\ge3$ portion of Theorem~\ref{thm:identity}. The
formula therefore expresses the tail of $Z_{\mathrm{1-loop}}$ in torsion-theoretic terms, and
does not reach $n=2$.
\end{proposition}

\begin{remark}[why the formula starts at $m\ge3$]
$\Rz(s,\sigma_k)$ converges absolutely for $\mathrm{Re}(s)>2$. The factor $n=2$ therefore sits
\emph{at} the abscissa of absolute convergence, not inside it. This is the structural reason the
ratio formula begins at $m\ge3$ and normalises by $\rho(2)$: the excluded factor is exactly the
one at the boundary of the region where the Euler product converges absolutely. It is also why
naive numerical evaluation of the $n=2$ term is oscillatory in the geodesic cutoff, while the
$n\ge3$ tail is not.
\end{remark}

\section{The $n=2$ factor at the abscissa}\label{sec:n2}

Because $n=2$ is the one factor outside the reach of Proposition~\ref{prop:tail}, its status
matters. We report a numerical study, and we are explicit that it is numerical.

\begin{proposition}[no breakdown, on an instrument that can see one]\label{prop:n2}
Approach the abscissa from above, where $\Rz(s,\sigma_2)=\prod_\gamma(1-e^{2i\theta_\gamma}
e^{-s\ell_\gamma})$ converges absolutely, and track the dependence on the geodesic cutoff. For
$m004$, over three cutoffs and seven values of $s$, $|\Rz(s,\sigma_2)|$ is smooth and monotone
through $s=2$ --- rising from $1.1075$ at $s=2.6$ to $1.1936$ at $s=2.0$ --- with no pole and no
discontinuity, and with a cutoff spread that grows smoothly rather than diverging.
\end{proposition}

\begin{scope}[what this does and does not establish]\label{sc:n2}
The reading is meaningful only because the instrument is shown able to report failure: taken
\emph{below} the abscissa, where divergence is certain, the same computation does diverge and
the cutoff spread blows up. That is the bite control, and without it a smooth curve would be
evidence of nothing.

What is established is the absence of any observable breakdown at $s=2$. What is \emph{not}
established is convergence. Three cutoffs are used and the cutoff-to-infinity limit is untested;
a spread that is small at cutoff $5.5$ may still fail to converge. Accordingly the status of the
$n=2$ factor moves from ``may not exist'' to ``no evidence of breakdown; existence unproved'',
and it remains one of the three residues below.
\end{scope}

\begin{remark}[a consistency check on the assembly]
Computing $\log Z$ for the assembled $n\ge3$ tail in two independent ways --- once through the
Ruelle factors $\Rz(k,\sigma_k)$ and once directly through the geodesic sum over $[\gamma]$ ---
agrees to $8.2\times10^{-14}$. This tests the identity of Theorem~\ref{thm:identity} in the
regime where both sides converge absolutely, rather than testing arithmetic.
\end{remark}

\section{The three residues}\label{sec:residues}

A single caveat covering everything unfinished is not useful to a reader. We separate the
obstructions, because they have different characters and different remedies, and because two of
the three are not what one would first guess.

\begin{residue}[the spin-$2$ cusped test function]\label{res:one}
Assembling the one-loop determinant for a finite-volume cusped hyperbolic $3$-manifold requires
a gauge-fixed spin-$2$ (and vector, and scalar) one-loop determinant in the cusped setting. We
have not found this in the literature, and we believe it to be a genuine gap rather than a
retrieval failure: adjacent literature returns richly for the compact case and for the
$\mathrm{SL}(2,\C)$ and $\mathrm{SL}(n,\C)$ state-integral families, while this specific
combination does not appear. \textbf{What would close it:} a construction and evaluation of that
determinant, or a citation we have missed.
\end{residue}

\begin{residue}[the evaluation point]\label{res:two}
Fried's theorem gives, for acyclic orthogonal representations, that the twisted Ruelle zeta is
holomorphic at $0$ with $\Rz_\sigma(0)=T_X(\sigma)^2$, and a cusped extension exists in the
literature via a regularised trace. It is natural to expect this closes the gap between the
graviton determinant and analytic torsion. \textbf{It does not, and the reason is the evaluation
point.} Fried's identity lives at $s=0$; Pfaff's ratios live at $k\ge3$; the graviton product
lives at $s=n\ge2$. These are three evaluations of one twisted Ruelle family, and Fried's is the
single point the graviton product never visits.

\textbf{The gap can nevertheless be located exactly.} Since $\rho(m)=\Sym^{2m}\C^2$ has
eigenvalues $e^{jL}$ for $j=-m,\dots,m$ on a holonomy of complex length $L=\ell+i\theta$, and
$e^{jL}e^{-s\ell}=e^{ij\theta}e^{-(s-j)\ell}$, one has the identity
\begin{equation}\label{eq:symfactor}
   \Rz_{\rho(m)}(s)=\prod_{j=-m}^{m}\Rz(s-j,\sigma_j),
\end{equation}
which is elementary and we claim no novelty for it. At $s=0$ it reads
\[
   \Rz_{\rho(m)}(0)=\Rz(0,\sigma_0)\cdot
   \underbrace{\prod_{j=1}^{m}\overline{\Rz(j,\sigma_j)}}_{\text{the graviton's own factors}}\cdot
   \underbrace{\prod_{j=1}^{m}\Rz(-j,\sigma_j)}_{\text{negative arguments}} .
\]
\textbf{The graviton's factors sit inside Fried's point explicitly}; the obstruction is the second
product, at arguments outside absolute convergence. \textbf{Residue~2 therefore reduces to
evaluating $\Rz$ at the negative integers} --- and since a reflection $s\leftrightarrow2-s$ carries
$2+j$ to $-j$, with $\Rz(2+j,\sigma_j)$ absolutely convergent for $j\ge1$, it reduces again:

\textbf{What would close it:} a functional equation for $\Rz(s,\sigma_k)$ on this manifold. That is
a single standard object rather than a wish across three theorems, and the relevant scattering
determinant is already available in closed form (\S\ref{sec:Lchi}). \textbf{We do not establish it
here, and residue~2 remains open.} We also do not verify Fried's hypotheses for $\rho(m)$; without
them, \eqref{eq:symfactor} still holds, but its composition with Fried's theorem does not follow.
\end{residue}

\begin{residue}[the cusp's continuous spectrum]\label{res:three}
The continuous spectrum contributed by the cusp is not included in the assembly above. Its
scattering determinant for $m004$ is available in closed form, but incorporating it into the
one-loop object is not done here. \textbf{What would close it:} the assembly, once
Residue~\ref{res:one} is available.
\end{residue}

\begin{remark}[a correction we record]
We initially expected Fried's theorem to close Residue~\ref{res:two} outright, and said so
before checking where each theorem evaluates. It does not. We record the expectation and its
failure because the reframing in Residue~\ref{res:two} is the useful output, and it was only
reached by asking which point of the family each theorem actually speaks about.
\end{remark}

\section{What the Ruelle zeta does not know about $L(s,\chi_{-3})$}\label{sec:Lchi}

The manifold $m004$ has invariant trace field $\Q(\sqrt{-3})$, and the Dirichlet $L$-function
$L(s,\chi_{-3})$ appears in its volume. It is natural to ask whether the Ruelle zeta factors
through that $L$-function. It does not.

\begin{theorem}\label{thm:Lchi}
The Euler product of $\Rz(k,\sigma_k)$ is purely geodesic: by Theorem~\ref{thm:identity} it is
$\prod_\gamma(1-q_\gamma^{\,k})$, indexed by primitive closed geodesics, and it contains no
$L$-function factor. In particular $\Rz$ does not factor through $L(s,\chi_{-3})$ in its Euler
product. However, $L(s,\chi_{-3})$ \emph{does} enter the functional equation, through the
scattering determinant of the cusp.
\end{theorem}

The distinction is worth preserving. Arithmetic enters this object through the cusp and the
functional equation, not through the geodesic Euler product; a claim that the Ruelle zeta ``is''
an $L$-function, or a ratio of them, is false in the Euler-product sense in which such claims
are usually meant.

\section{Independent corroboration}\label{sec:corrob}

The two central negative findings of this paper --- that the cited torsion formulae and the
scalar cusp scattering determinant do not by themselves supply the one-loop object, and that a
gauge-fixed spin-$2$ determinant for the finite-volume cusped case is not available --- were
reached independently, by different means, and recorded in a register of questions as
\texttt{REFUTED} and \texttt{EXTERNAL\_BLOCKER} respectively (Appendix~\ref{app:register}).

We note this because the second is an assertion of absence, and assertions of absence are the
weakest kind of claim a single investigator can make. Two independent arrivals at the same
boundary --- agreeing not only that something is missing but on \emph{which} thing is missing and
which adjacent things are present --- is materially better evidence than one, and it is the
reason we are willing to state Residue~\ref{res:one} as a gap rather than as a failure to find.

\section{What is not claimed}

No measured quantity appears in this paper, and no identification of any quantity here with any
observable is offered or used. The numerical work of \S\ref{sec:n2} is evidence about
convergence behaviour, not a computation of a physical value. Convergence of the $n=2$ factor is
unproved (Scope~\ref{sc:n2}), and the one-loop partition function is not assembled
(\S\ref{sec:residues}).

\appendix

\section{Verification}\label{app:verify}

The computations are reproduced by the scripts in \texttt{verify/}, over the complex length
spectrum of $m004$. Controls run before results are read.

\begin{center}
\begin{tabular}{@{}lp{8.4cm}@{}}
\toprule
discharges & control or check \\
\midrule
Thm~\ref{thm:identity} & $\log Z$ for the $n\ge3$ tail computed twice, once via $\Rz(k,\sigma_k)$ and once via the direct geodesic sum, agreeing to $8.2\times10^{-14}$; conjugate pairs identified by distance to the nearest multiple of $2\pi$, not by a raw modulus, and the four self-conjugate $\theta=\pi$ classes handled separately \\
Prop.~\ref{prop:tail}  & the torsion ratios at $m=3,4,5,6$ reproduce independently computed values; $\kappa(X)=1$; $\vol(X)=2.029883212819307$ taken at full double precision, since a truncated decimal string manufactures a spurious mismatch at the $10^{-11}$ level \\
Prop.~\ref{prop:n2}    & \textbf{bite control:} the same evaluation taken \emph{below} the abscissa must diverge with a blowing-up cutoff spread, and does; the reported smoothness above the abscissa is meaningful only against it \\
Thm~\ref{thm:Lchi}     & the Euler product is regenerated from the geodesic list rather than quoted, and checked against a Dirichlet-series evaluation performed by analytic continuation --- a naive series at $s=1/2$ lies outside its half-plane and produces a spurious refutation \\
\bottomrule
\end{tabular}
\end{center}

\section{Status in the programme register}\label{app:register}

\begin{center}
\begin{tabular}{@{}llp{6.9cm}@{}}
\toprule
entry & status & relation to this paper \\
\midrule
\texttt{OA-C1061} & \texttt{REFUTED} & whether the cited Fried, Park or Pfaff torsion formulae, or the scalar cusp scattering determinant, directly supply the one-loop object. Independently reached; matches Corollary~\ref{cor:notasingle} and Residue~\ref{res:two}. \\
\texttt{OA-C1062} & \texttt{EXTERNAL\_BLOCKER} & whether a gauge-fixed spin-$2$/vector/scalar one-loop determinant can be constructed for the finite-volume cusped case. Independently reached; this is Residue~\ref{res:one}. \\
\texttt{OA-C1059} & \texttt{REFUTED} & whether the untwisted Ruelle zeta is a finite product or ratio of shifted Dirichlet $L$-functions. Consistent with Theorem~\ref{thm:Lchi}. \\
\texttt{OA-C1060} & \texttt{EXTERNAL\_BLOCKER} & whether the $n=2$ factor has a proved cutoff-independent value. This is exactly Scope~\ref{sc:n2}: we report no breakdown and do not claim existence. \\
\bottomrule
\end{tabular}
\end{center}

\noindent
Three of the four are negative or blocked. We regard that as the honest shape of this subject at
present, and prefer to publish the identity together with an accurate map of what it does not
close.

\begin{thebibliography}{9}
\bibitem{Fried} D.~Fried, Analytic torsion and closed geodesics on hyperbolic manifolds,
  \emph{Invent.\ Math.} \textbf{84} (1986), 523--540.
\bibitem{Pfaff} J.~Pfaff, Analytic torsion versus Reidemeister torsion on hyperbolic
  $3$-manifolds with cusps, arXiv:1206.0228.
\bibitem{Park} J.~Park, Analytic torsion and Ruelle zeta functions for hyperbolic manifolds
  with cusps, \emph{J.\ Funct.\ Anal.} \textbf{257} (2009), 1713--1758.
\bibitem{GMY} S.~Giombi, A.~Maloney and X.~Yin, One-loop partition functions of $3d$ gravity,
  \emph{JHEP} \textbf{08} (2008), 007.
\bibitem{MaloneyWitten} A.~Maloney and E.~Witten, Quantum gravity partition functions in three
  dimensions, \emph{JHEP} \textbf{02} (2010), 029.
\bibitem{Melrose} R.~B.~Melrose, \emph{The Atiyah--Patodi--Singer Index Theorem}, A K Peters,
  1993.
\end{thebibliography}

\end{document}
